\documentclass[UTF8]{ctexbook}
\usepackage{ctex}
\usepackage{enumerate}
\usepackage{hyperref}
\pagestyle{plain}
\title{药理学}
\author{Sakura}

\begin{document}
\tableofcontents
\setcounter{chapter}{22}
\chapter{肾素—血管紧张素系统药理}
了解肾素-血管紧张素系统；掌握ACE抑制药的药理作用与临床应用及不良反应。
\section{ACE抑制药}
\subsubsection{化学结构和分类}
ACE有两个结合位点，其中含Zn2+的是必须结合的位点，一旦结合ACE活性消失。一共有三类

1、含有巯基：卡托普利；2、含有羧基：依那普利、雷米普利、培哚普利、贝那普利、赖诺普利；3、含有磷酸基的：福辛普利。

许多ACE抑制药是前体药物

\subsubsection{药理作用}
1、阻止AngI转化为 Ang II，取消AngII的收缩血管、刺激醛固酮释放、增加血容量、升高血压和促心血管肥大增生的作用，有利于高血压、心力衰竭和心血管重构的防治

2、保存缓激肽活性。缓激肽可以激活缓激肽B2受体，使NO和PGI2生成增加，舒张血管、降低血压、抗血小板聚集、抗心血管细胞肥大增生。

3、保护血管内皮细胞

4、抗心肌缺血与心肌保护，减轻心肌缺血再灌注损伤

5、增敏胰岛素受体
\subsubsection{临床应用}
1、治疗高血压，可联用利尿药增效。对肾血管性高血压因肾素水平高特别有效。

2、治疗充血性心力衰竭或心肌梗死。

3、治疗糖尿病性肾病和其他肾病。能够阻止肾小囊内压升高导致的肾小球和肾功能损伤。与降压作用无关，为舒张肾小球出球小动脉的结果
\subsubsection{不良反应}
1、首剂低血压。口服吸收快、生物利用度高的多见；

2、咳嗽，因为缓激肽或前列腺素、P物质在肺内蓄积的结果。

3、高血钾，使依赖AngII的醛固酮分泌减少导致血钾增高，如同时服用保钾利尿药患者常见

4、低血糖

5、肾功能损伤。ACE抑制药舒张出球小动脉，降低肾灌注压，导致肾滤过率和肾功能降低。

6、妊娠第二三期致畸

7、血管神经性水肿

8、含—SH的ACE可以导致味觉障碍、皮疹和白细胞缺乏等和其他—SH药物相似的反应。
\chapter{利尿药}
了解尿的形成过程；掌握脱水药的药理作用和作用机制；重点掌握利尿药分类、作用机制、临床应用及不良反应。
\section{碳酸酐酶抑制药}作用于近曲小管。
\subsection{乙酰唑胺}
\subsubsection{药理作用}
通过抑制碳酸酐酶的活性而抑制HCO3-的重吸收，治疗量的乙酰唑胺可以抑制近85\%的重吸收，近曲小管Na+重吸收会减少，水的重吸收减少，但是集合管Na+重吸收大大增多，使K+的分泌增多。
\subsubsection{临床应用}
1、治疗青光眼，减少房水生成降低眼内压；2、机型高山病，可减少脑脊液的生成和pH，减轻症状改善机体功能；3、碱化尿液，可促进弱酸性物质排泄，但只在第一次有效，长期服用需要补充碳酸氢盐；4、纠正代谢性碱中毒；5、其他：癫痫辅助治疗、伴随低钾血症的周期性麻痹、严重高磷酸盐血症。
\subsubsection{不良反应}
1、过敏，磺胺衍生物；2、代谢性酸中毒；3、尿结石；4、失钾；5、嗜睡感觉异常等其他毒性。
\section{渗透性利尿药}脱水药。作用于髓袢和肾小管其他部位。甘露醇、山梨醇和高渗葡萄糖。脱水、利尿作用。可预防急性肾衰竭。
\section{袢利尿药}高效能利尿药或\(\textup{Na}^{+}-\textup{K}^{+}-\textup{2Cl}^{-}\)同向转运子抑制药，作用于髓袢升支粗段，利尿作用强、代表药物呋塞米、依他尼酸、布美他尼。
\subsubsection{药理作用}特异性的与Cl-结合位点结合而抑制分布在髓袢升支管腔膜侧的\(\textup{Na}^{+}-\textup{K}^{+}-\textup{2Cl}^{-}\)同向转运子，抑制NaCl的重新收。由于K+重吸收减少，降低了K+的再循环，导致管腔正电位、减少Ca2+和Mg2+的重吸收，排泄增加。在远曲小管和集合管位，\(\textup{Na}^{+}-\textup{K}^{+}\)交换增多，进一步促进K+排出。

袢利尿药可以促进肾脏前列腺素的合成，因此非甾体抗炎药可以干扰利尿药的作用

能够调节血管，对心力衰竭的患者能够在利尿作用发生前就产生有效的血管扩张作用。呋塞米和依他尼酸可以迅速增加全身静脉血容量，降低左室充盈压，减轻肺淤血
\subsubsection{临床应用}
1、急性肺水肿和脑水肿，可以迅速扩张血管，在利尿作用发生之前即可缓解急性肺水肿。因为血液浓缩对脑水肿也有效。

2、其他严重水肿。

3、急慢性肾功能衰竭。急性肾衰时，可以增加尿量和K+的排出，冲洗肾小管，减少萎缩和坏死，但不能延缓肾衰的进程。大剂量呋塞米可以治疗慢性肾衰，增加尿量，其他药物无效时仍然能产生作用。可以扩张肾血管、增加肾血流量和肾小球滤过率。

4、高钙血症。抑制Ca2+排泄，可联合静脉注射生理盐水加快。

5、加速某些毒物排泄。
\subsubsection{不良反应}
1、水和电解质紊乱，低血容量、低血钾、低血钠、低氯性碱血症、长期可致低血镁；2、耳毒性，和内耳淋巴液电解质成分改变有关；3、高尿酸血症：血容量降低、细胞外液容积减少，尿酸经近曲小管重吸收增加，也和竞争有机酸分泌途径有关；4、其他：高血糖、恶心呕吐、胃肠道出血、升高LDL胆固醇甘油三酯、降低HDL
\section{噻嗪类利尿药和类噻嗪类利尿药}中效能利尿药\(\textup{Na}^{+}-\textup{Cl}^{-}\)同向转运子抑制药，作用于远曲小管近端

氯噻吩（噻吩）；吲达帕胺、氯噻酮、美托拉宗、喹乙宗（磺胺）。
\subsubsection{药理作用}
利尿作用，增加NaCl和水的排出，K+排出也增多，依赖于前列腺素的产生并受到解热镇痛药的抑制

抗利尿作用：减少尿崩症患者的尿量和口渴症状。

降压作用，用药早期用过利尿、减少血容量，长期可以扩张外周血管。
\subsubsection{临床应用}
1、水肿；2、高血压；3、其他：肾性尿崩症和加压素无效的垂体性尿崩症。
\subsubsection{不良反应}
1、电解质紊乱；2、高尿酸血症（竞争有机酸通道）；3、代谢变化：高血糖高血脂，增加LDL、胆固醇、抑制胰岛素分泌；4、过敏反应，本品为磺胺类可以交叉过敏。
\section{保钾利尿药}低效能利尿药，作用于远曲小管远端和集合管。
分为两类，醛固酮受体拮抗剂和肾小管上皮细胞Na+通到道抑制药。
\subsection{螺内酯}
\subsubsection{药理作用}
螺内酯及其代谢产物坎利酮与醛固酮相似，能够结合到远曲小管细胞胞质中的盐皮质受体，阻止醛固酮-受体复合物的核转位，而产生拮抗醛固酮的作用。
\subsubsection{临床应用}
保钠排钾作用，利尿作用弱。主要用于醛固酮升高有关的顽固性水肿和充血性心力衰竭。
\subsection{氨苯喋呤和阿米洛利}
作用于远曲小管末端和集合管，阻滞管腔Na+通到而减少Na+的重吸收。负电位降低因此K+分泌驱动力弱。并非拮抗醛固酮，对肾上腺切除患者仍有作用。主要用于治疗顽固性水肿。
\chapter{抗高血压药}
了解血压的形成及调控，高血压的分期及病理改变；重点掌握抗高血压药的分类及作用特点。
\section{利尿药}噻嗪类常用。
\section{钙通道阻滞药}
减少细胞内钙离子含量松弛血管平滑肌，降低血压
\subsection{硝苯地平}
\subsubsection{药理作用}
1、作用于细胞膜L型钙通道，抑制钙离子从细胞外进入细胞内，是小动脉扩张，总外周血管阻力下降。冠状动脉和脑血管较为敏感，对子宫平滑机送出作用较为明显。

2、负性肌力作用，心肌收缩力降低，耗氧量减少。负性频率和负性传导作用作用，使0相除极和4相缓慢除极减慢，降低室房结传导速度，减慢心率。

3、抗动脉粥样硬化作用，钙离子参与动脉粥样硬化的病理过程，钙通道阻滞药可以干扰。（1）减少Ca2+超载的动脉壁损坏；（2）抑制平滑肌增殖和动脉基质蛋白质合成，增加血管顺应性；（3）抑制脂质过氧化，保护内皮细胞；（4）增加胞内cAMP含量，降低细胞内胆固醇水平。

4、减轻Ca2+超负荷对红细胞的损伤，地尔硫卓能抑制TXA2的产生和ADP、Adr和5-HT引起的血小板聚集。

5、保护肾脏，对伴随肾功能障碍的高血压病患者。
\subsubsection{临床应用}
轻中重度高血压。心绞痛，变异性心绞痛、稳定型心绞痛和不稳定型心绞痛。心律失常。脑血管疾病。
\section{\(\beta\)受体阻断药}
\subsection{普萘洛尔}
\subsubsection{药理作用}
非选择性\(\beta\)受体阻断药，对\(\beta_{1}\)和\(\beta_{2}\)有相同的亲和力，缺乏内在拟交感活性。可通过减少心输出量、抑制肾素释放、抑制交感神经系统活性、增加前列环素的合成等来产生降压作用。
心率减慢，心肌收缩力和心输出量降低，冠脉血流量减少，心肌耗氧量减少，对高血压患者可以降压
\subsubsection{临床应用}
心律失常、心绞痛、高血压、甲亢。
\section{ACEI抑制药}
\subsection{卡托普利}
抑制ACE，使AngI转化为AngII减少，舒张血管。同时减少醛固酮分泌以利于排钠。
\section{AT1阻断药}
\subsection{氯沙坦}
竞争性阻断AT1受体，对抗AngII，用于高血压。
\section{中枢性降压药}
\subsection{可乐定}
\subsubsection{药理作用}
兴奋延髓背侧孤束核突触后膜的\(\alpha_{2}\)受体，抑制交感神经的活性，使外周血管扩张，血压下降。
\subsubsection{临床应用}
中重度高血压，常用语其他药无效时，降压作用中等偏强，不影响肾血流量和肾小球滤过率，可长期用于高血压治疗。和利尿药合用可用于重度高血压。也可用于预防偏头痛和吗啡成瘾戒毒药。滴眼液用于青光眼。
\subsubsection{不良反应}
口干、便秘。嗜睡、抑郁、眩晕、血管性水肿、腮腺肿痛、心动过缓、食欲不振。不适合用于高空作业患者和驾驶机动车辆的人员。

\chapter{治疗心力衰竭的药物}
了解CHF的病理生理学；掌握治疗CHF药物的分类；重点掌握肾素-血管紧张素-醛固酮系统抑制药治疗CHF的作用机制，强心苷类的药理作用、作用机制、临床用途与不良反应。
\section{利尿药}促进Na+的排泄和水的排泄，血容量减少，降低心脏前负荷，改善新功能；降低静脉压，消除或缓解静脉淤血及其一发的肺水肿和外周水肿。
\section{肾素-血管紧张素-醛固酮系统抑制药}
\subsection{ACE抑制药}
1、降低外周血管阻力，降低心脏后负荷；2、减少醛固酮生成，减轻水钠潴留，降低心脏前负荷；3、抑制心肌及血管重构；4、降低全身血液阻力，能降低左心室充盈压、左室舒张末压，降低室壁张力，改善心脏舒张功能，降低肾血管阻力增加肾血流量；5、降低交感神经活性。
\subsection{AT1受体拮抗药}
拮抗AngII的作用，具有ACE抑制药的所有好处但减少不良反应。
\subsection{抗醛固酮药}
抑制心肌与血管重构，防止左室肥厚时心肌间质纤维化。
\section{\(\beta\)受体阻断药}
\subsubsection{作用机制}
1 拮抗交感活性：

1) 阻断心脏\(\beta\)受体、拮抗过量儿茶酚胺对心脏的毒性作用，防止大量Ca2+内流所诱发的心肌损伤；

2) 改善心肌重构

3）减少肾素释放，抑制RAS，防止高浓度AngⅡ对心脏的损害；

4）上调心肌\(\beta\)R，改善\(\beta\)R对儿茶酚胺的敏感性(卡维地洛除外）

5）卡维地洛兼阻断\(\alpha_{1}\)受体、抗氧化等作用。

2、抗心律失常和心肌缺血的作用。
\section{正性肌力药物}
\subsection{强心苷}
\subsubsection{化学组成}
由苷元和糖组成，苷元是强心苷作用的关键部位。

1）C3位\(\beta\)-OH是甾核与糖的结合部位，若脱掉糖，转为\(\alpha\)构型失效；2）C14位\(\beta\)-羟基（强心必需)；3）C17位连接\(\beta\)构型的不饱和内酯环（打开，被饱和失效）4） 甾核的羟基数目多，则作用快，持续时间短

毒毛花苷-K（4个羟基） 速效、短效

洋地黄毒苷（1个羟基） 慢效、长效
\subsubsection{药理作用}
1、正性肌力作用。能显著增加衰竭心脏的心肌收缩力，增加心输出量，从而解除心衰的症状。\textcircled{1}加快心肌纤维缩短速度，是心肌收缩敏捷，舒张期相对延长；\textcircled{2}加强衰竭心肌收缩力，增加心输出量的同时不增加耗氧量，甚至有所降低。

2、负性心率作用。对心率加快及伴有房颤的新功能不全者可显著减慢心率。新功能不全时由于反射性交感神经活性增加，是心率加快。使用强心苷后心输出量增加，反射性的兴奋迷走神经，使窦房结抑制，心率降低。

对神经和内分泌系统的作用：1、中毒剂量强心苷兴奋延髓催吐化学感受区－－呕吐，并使交感神经冲动发放增加，引起快速型心律失常。2、强心苷抑制心功不全时过度激活的 RAAS。

利尿作用：1、心功能改善后，增加肾血流量和肾小球滤过功能 。2、直接抑制Na+-K+-ATP酶，肾小管重吸收Na+减少，钠和水排出增多

对血管的作用：正常人：能接收缩血管平滑肌，使外周阻力上升；CHF患者用药后，交感神经活性减低的作用超过直接收缩血管的效应。故外周阻力减小，心排出量增加，局部血流量增加。
\subsubsection{临床应用}
1、治疗心力衰竭。

1）对伴有心房纤颤或心率较快的CHF疗效最佳。

2）对瓣膜病、风湿性心脏病（高度二尖瓣狭窄的病例除外）、冠状动脉硬化性心脏病和高血压性心脏病所致的CHF疗效较佳。

3）对机械性阻塞和能量代谢障碍的CHF疗效差。

4）对肺源性心脏病、活动性心肌炎疗效较差。

5）对扩张型心肌病、舒张性CHF---不用

2、某些心律失常。

心房纤颤：心房肌发生细弱而不规则纤维颤动。危害：心率过快，严重循环障碍。治疗：兴奋迷走神经或对房室结的直接作用，减慢房室传导。减慢心室率，改善循环障碍。但多数病人不能终止。

心房扑动，最常用药物。兴奋迷走神经，促进K+外流，心房肌细胞RP加大、提高0相除极速率,心房的传导上升。故缩短心房ERP，扑动变为颤动，在心房纤颤时更易增加房室结隐匿性传导而减慢心室率。部分病人在纤颤后停用强心苷可恢复窦性节律。

阵发性室上性心动过速，增强迷走神经兴奋性，降低心房兴奋性。

\subsubsection{不良反应}
1、心脏反应：

1）快速型心律失常：最多见的最早见的是室性前期收缩。与强心苷抑制\(\textup{Na}^{+}-\textup{K}^{+}-\textup{ATP}\)酶被高度抑制外，也与强心苷引起的迟后除极有关。

2）房室传导阻滞：与提高迷走神经兴奋性有关，也有抑制\(\textup{Na}^{+}-\textup{K}^{+}-\textup{ATP}\)有关。

3）都行心动过缓：强心苷可因为抑制窦房结、降低其自律性而引起心动过缓，可作为停药指征。

氯化钾是治疗强心苷中毒引起的快速型心律失常的有效药物，能与强心苷竞争结合\(\textup{Na}^{+}-\textup{K}^{+}-\textup{ATP}\)酶。但是钾离子结合较为疏松，因此只能阻止强心苷继续结合而不能置换已经结合的心肌细胞。因此注意防止低血钾更重要。钾离子不可过量，还要注意肾功能防止高血钾。并发传导阻滞的不能补钾盐。心律失常严重者还应使用苯妥英钠，苯妥英钠可以抗心律失常，还可以将与\(\textup{Na}^{+}-\textup{K}^{+}-\textup{ATP}\)酶结合的强心苷置换出来。

利多卡因可用于治疗强心苷引起的室性心动过速和心室纤颤。心动过缓和房室传导阻滞不能补钾盐，可用M受体阻断剂如阿托品。

2、胃肠道反应：恶心、厌食、呕吐等。

3、中枢神经系统反应：炫目、头晕、失眠、疲倦。视觉异常通常是中毒先兆可作为停药指征。

\subsection{非苷类正性肌力药}
可能增加心衰患者的病死率 不宜用作常规治疗用药

儿茶酚胺类－－用于强心苷反应不佳的CHF

多巴酚丁胺：\(\beta_{1}\)受体激动剂增强心肌收缩性，心输出量（但不明显改变心率），用于强心苷反应不佳的CHF，降低外周阻力。

磷酸二酯酶抑制药（phosphodiesteraseinhibitor）氨力农amrinone和米力农milrinone：抑制磷酸二酯酶Ⅲ，心肌细胞内cAMP含量上升，cAMP激活PKA，促进Ca2+内流，正性肌力和血管舒张，心输出量增加，心脏负荷减小，耗氧量减小。用于心衰时作短时间的支持疗法。1）氨力农主要用于治疗严重及对强心苷和利尿药不敏感的CHF；2）氨力农不良反应较重：恶心，呕吐，心律失常发生率较高等；3）米力农为氨力农的替代品，不良反应较少。
\section{扩血管药}
1、扩张静脉，静脉回心血量减少，心脏前负荷降低 ，肺楔压下降 、左心室舒张末压下降 ，缓解肺部淤血。

2、扩张小动脉，外周阻力下降 ，心脏的后负荷降低 ，增加心输出量，增加动脉供血，缓解组织缺血。

3、 辅助用药：1） 肺静脉压明显上升，肺淤血明显——扩张静脉为主的硝酸酯；2）外周阻力上升 ，心输出量下降——扩张小动脉为主的肼屈嗪（适应症：肾功能不良或不耐ACEI患者）；3）肺静脉压和外周阻力上升，心输出量下降，肺淤血——扩张动静脉如哌唑嗪
\section{钙通道阻滞药及钙增敏药}
钙通道阻滞药：1、扩张血管 降低心脏前后负荷；2、激活交感神经系统和负性肌力作用；3、钙拮抗药对收缩性心室功能障碍者并不降低病死率；4、适应症－－ 舒张功能障碍的心衰

钙增敏药：作用于收缩蛋白，增加肌钙蛋白C对钙离子的亲和力，在不增加细胞内钙浓度下，增强心肌收缩能力。----疗效有待于进一步确证。
\chapter{调血脂药与抗动脉粥样硬化药}
了解血浆脂蛋白的分类，调血脂药物的临床应用与不良反应；掌握调血脂药物的分类和作用。

脂类包含胆固醇(Ch)、三酰甘油(TG)、磷脂(PL)、游离脂肪酸(FFA)，Ch又分为游离(FC)和胆固醇酯(CE)，两者相加为总胆固醇(TC)。血脂和载脂蛋白Apo结合形成脂蛋白LP后才能溶于血浆、进行转运和代谢。分为乳糜微粒CM、极低密度脂蛋白VLDL、低密度脂蛋白LDL、中间密度脂蛋白IDL和高密度脂蛋白HDL。IDL是VLDL在血浆中的代谢物。
乳糜微粒(CM)，主要为TG, 占85\%，功能：转运外源性TG和CH

极低密度脂蛋白VLDL， 55\%TG,功能：转运内源性TG到周围组织

LDL ，主成分：胆固醇酯（37\%）和ApoB，功能：将内源性CH转运到周围组织细胞

HDL，主成分：较多的ApoA, 少量ApoC，Ch和CE，功能：外周组织的CH转运到肝分解代谢，即逆向转运。

高脂蛋白血症，血浆中LDL和VLDL 增多。原发性：家族性，继发性：常见糖尿病、肾病综合征、慢性肝病等。血浆中VLDL、IDL、LDL及apo B浓度高出正常为高脂蛋白血症，易致动脉粥样硬化。HDL、apo A浓度低于正常，也为动脉粥样硬化危险因子。
\section{调血脂药}
主要降低TC和LDL的药物：包括他汀类、胆汁酸结合树脂、酰基辅酶A胆固醇酰基转移酶抑制药。

主要降低TG及VLDL的药物：贝特类药物和烟酸。

降低 LP（a）的药物：现已证明烟酸、烟酸戊四醇酯、烟酸生育酚酯、阿西莫司等。
\subsection{他汀类}HMD-CoA还原酶抑制剂，该酶是肝细胞合成胆固醇的限速酶。可以减少胆固醇的合成。他汀类具有二羟基庚酸结构或为内酯或为开环羟基酸，是抑制HMG-CoA还原酶的必需基团
\subsubsection{药理作用}
1、调血脂作用及作用机制：对LDL-C降低作用最强，TC次之，降TG作用弱。剂量依赖性，两周起效，4-6周高峰，HDL-C略升高，长期可保持疗效。他汀类药物与HMG-CoA相似，可以竞争性抑制还原酶，使胆固醇合成受阻。代偿性增加肝细胞表面LDL R，使血浆LDL降低，VLDL代谢加快，肝合成和释放VLDL降低。

2、非调血脂性作用，但有助于抗动脉粥样硬化：

1）改善血管内皮功能，提高血管内皮对扩血管物质的反应性

2）抑制血管平滑肌细胞VSMCs增殖和迁移，促进VSMCs凋亡

3）降低血浆C反应蛋白，炎性反应减轻

4）抑制单核细胞-巨噬细胞的黏附和分泌功能

5）抑制血小板聚集和提高纤溶活性等 。

6) 抗氧化，清除自由基，抑制ox-LDL的生成

7）动脉壁巨噬细胞及泡沫细胞形成减少，动脉斑块稳定和缩小。基质金属蛋白酶（MMP）：分解基质，加速胶原降解，降低纤维帽的抗张强度，斑块破裂。TNF损伤结构蛋白，增加纤维帽的脆性；刺激细胞表达MMP，使斑块易于破裂。他汀下调MMP表达，降低巨噬细胞活性，降低斑块T细胞活性，下调TNF含量，产生抗AS作用。

3、肾脏保护：纠正脂质代谢异常而引起的慢性损害。
\subsubsection{临床应用}
1、调节血脂：杂合子家族性和非家族性\(\textup{II}_{a}\)、\(\textup{II}_{b}\)、\(\textup{III}\)型高脂蛋白血症。也可用于糖尿病和肾病综合征引起的高胆固醇血症。

2、肾病综合征：除与调血脂作用有关外，还与抑制肾小球膜细胞增殖，延缓肾动脉硬化有关。

3、预防心脑血管急性时间：可以加强粥样斑块的稳定性或使斑块缩小。

4、抑制血管成形术后再狭窄。
\subsubsection{不良反应}少
\subsection{胆汁酸结合树脂}与胆汁酸结合阻止肝肠循环，大量消耗胆固醇，使TC和LDL-C水平降低。
考来烯胺和考来替泊
\subsubsection{药理作用}
1、能减少食物中之类的吸收；2、阻滞胆汁酸在肠道的重吸收；3、由于大量丢失胆汁酸，肝脏大量转换胆固醇为胆汁酸；4、由于肝内胆固醇减少，LDL受体增加或活性增强；5、LDL-Ch经受体进入肝细胞，因此血浆TC和LDL-Ch减少；6、HMG-CoA有继发性活动增强但是不能抵消，可联合使用他汀类药物。
\subsubsection{临床应用}
适用于\(\textup{II}_{a}\)、\(\textup{II}_{b}\)和家族性杂合子高脂血蛋白症，对纯合子家族性高胆固醇血蛋白症无效，对\(\textup{II}_{b}\)应与降TG和VLDL的药物配合使用。
\subsubsection{不良反应}
便秘、腹胀、嗳气、食欲减退。长期应用影响脂溶性维生素，叶酸及其他一些药物吸收
\subsection{酰基辅酶A胆固醇酰基转移酶(ACAT)抑制药}甲亚油酰胺
1、抑制ACAT，阻滞细胞内Ch向胆固醇酯（CE）的转化，减少外源性Ch的吸收，肝VLDL形成减少

2、阻滞外周组织CE的蓄积和泡沫细胞的形成，有利于Ch的逆化转运。

3、 抑制ACAT，发挥调血脂和抗ＡＳ的效应

4、适用于Ⅱ型高脂蛋白血症

5、 不良反应轻微， 可有食欲减退或腹泻。
\subsection{贝特类}
\subsubsection{药理作用}
既有调血脂作用也有非调血脂作用，能降低血浆TG、VLDL-C、TC、LDL-C，能升高HDL-C。各种不同药物作用强度不同。非调血脂作用有抗凝血、抗血栓和抗炎作用等。

激活类固醇激素受体类的核受体-过氧化物酶体增殖受体\(\alpha\)，调节LPL、ApoCIII、ApoI等基因的表达，降低ApoCIII转录、增加LPL和Apo A I的生成和活化有关。
\subsubsection{临床应用}原发性高TG血症，对III型高脂蛋白血症和混合型高脂蛋白血症和II型糖尿病的高脂蛋白血症有效。
\section{其他抗动脉粥样硬化药物}
抗氧化剂：普罗布考、多烯脂肪酸、黏多糖和多糖类：低分子量肝素。
\chapter{抗心绞痛药}
了解心绞痛的概念、发病机制和临床分型

掌握抗心绞痛药的分类、作用机制和作用特点

重点掌握硝酸酯类、\(\beta\)阻断剂和钙拮抗剂治疗心绞痛的作用机理及应用

冠状动脉性心脏病（coronary heart disease），简称冠心病（缺血性心脏病）：冠状动脉粥样硬化或痉挛时所导致的心肌缺血缺氧。心绞痛是冠状动脉供血不足引起心肌急剧的、暂时的缺血和缺氧而引起的临床综合症。

临床表现：胸骨后部或左前胸出现阵发性绞痛闷痛或压榨性痛感（重点），常放射至左肩，左上肢及肢端。一般历时1~5min，很少超过10~15min。休息或用药物治疗后常可缓解。痛觉的因素：缺血缺氧时积聚过多的代谢产物，如乳酸、丙酮酸、K+，以及组胺，刺激心肌植物神经传人纤维末梢。
\section{分类}
劳累性心绞痛(angina of effort)：稳定型、初发型、恶化型由劳累、情绪波动或其他增加心肌耗氧量的因素引发，休息或舌下含服硝酸甘油可缓解。稳定型心绞痛：心外膜冠状动脉的粥样硬化性狭窄所致。

自发性心绞痛 (angina pectoris at rest)：与心肌耗氧量无明显关系，多发于安静状态，症状重时间长。不易被硝酸甘油缓解 。

混合性心绞痛 (mixed pattern of angina)：心肌需氧量增加或无明显增加时都有可能发生为劳累性和自发性心绞痛混合出现。

\section{概述}生理机制：心肌耗氧和心肌供氧平衡失调。

供氧：取决于动静脉的氧分压差及冠状动脉的血流量；增加氧供应：主要增加冠状动脉的血流量

心室壁张力（心室壁张力:与心室内压力和心室容积成正比，与心室壁厚度成反比。）、心肌收缩力、收缩速率、心率、心输出量与耗氧成正比。

\section{作用机理}.1
\subsubsection{降低心肌耗氧量}
1、扩张血管：静脉扩张，则回心血量下降，降低前负荷；动脉扩张，降低后负荷，降低心肌耗氧量

2、减慢心率、 降低心肌收缩力
\subsubsection{提高心肌供氧，增加冠脉血流量}
1、舒张冠状动脉，解除动脉痉挛；2、促进侧支循环的形成

\section{硝酸酯类}
\subsection{硝酸甘油；单硝酸异山梨酯；硝酸异山梨醇酯}
\(\textup{-O-NO}_{2}\)是关键结构。

口服效果差，首关效应大，主要用于舌下含服（一定程度避免首关效应），维持20-30分钟；

肝中代谢，被谷胱甘肽还原酶代谢为二硝酸代谢产物，肾脏排出

长效：硝酸异山梨醇酯（消心痛）, 维持45-60分钟(舌下)

\subsubsection{药理作用}
1、降低心肌耗氧量扩张静脉血管，回心血量减少，心室容积缩小，心室内压减小，心室壁张力降低，射血时间缩短，前负荷减小。
较大剂量可扩张动脉血管，降低心脏的射血阻力，从而降低左室内压和心脏后负荷而降低心肌耗氧量。
虽然会引起反射性的心率加快，但总体结果是耗氧量降低。

2、扩张冠状动脉 增加缺血区血液灌注扩张较大的心外膜血管、输送血管及侧支血管－－冠状动脉痉挛更为明显，对阻力血管舒张作用弱。
因此导致非缺血区阻力大于缺血区，血液顺输送血管从侧支流入缺血区

3、降低左室充盈压，增加心内膜供血，改善左室顺应性
（1）扩张静脉血管，减少回心血量，降低心室内压；（2）扩张动脉血管，降低心室壁张力，增加了心外膜向心内膜的有效灌注压，增加心内膜血流量

4、保护缺血的心肌细胞，减轻缺血损伤
（1）释放NO，促进内源性的PGI2，降钙素基因相关肽等物质生成与释放。——保护；（2）缩小心肌梗死范围，改善左室重构，减少心肌缺血合并症。
\subsubsection{作用机制}
1、作为NO的供体，在平滑肌细胞内经谷胱甘肽转移酶的催化释放出NO，NO激活GC，细胞内cGMP含量上升，激活蛋白激酶，减少细胞内Ca2+释放和外Ca2+内流，胞浆中Ca2+下降，促使肌球蛋白轻链（MLCK)去磷酸化而松弛血管平滑肌。

2、\(\textup{PGI}_{2}\)－－扩血管

3、硝酸甘油促进降钙素基因相关肽合成及释放－－后者激活ATP敏感型钾通道，细胞膜超级化，扩血管。

4、NO抑制血小板聚集和粘附，有利冠心病的治疗。
\subsubsection{临床应用}
可用于治疗各型心绞痛、进行心肌梗死（缩小梗死范围、抗血小板聚集和粘附）、可用于慢性心力衰竭（CHF）

\subsubsection{不良反应}
血管扩张引起的：面颈部皮肤发红、搏动性头痛、眼内压升高

大剂量：体位性低血压及晕厥

剂量过大：灌注压力过低，反射性引起心率上升和心肌收缩力增加、加重心绞痛

超剂量：高铁血红蛋白血症（呕吐和紫癜）

连续用药出现耐受性，停药1~2周后，耐受性可消失。
1、与反复用药引起血管平滑肌细胞内巯基耗竭，NO生成障碍有关。 －－血管耐受；2、硝酸酯类使血管内压力下降，反射性增强交感活性，激活肾素－血管紧张素系统等－－伪耐受。
可间歇给药或给予含巯基的药物。
\section{\(\beta\)肾上腺素受体阻断剂}
\subsection{普萘洛尔（心得安）}
不适于变异型心绞痛（冠状动脉痉挛诱发），心绞痛发作次数减少，心电图改善，BP减少。
\subsubsection{作用机制}
降低心肌耗氧量：阻断心肌\(\beta\)受体使心率下降，收缩力下降，耗氧量下降，但抑制心肌收缩力可增加心室容积，同时因收缩力下降，心室射血时间增长，心肌耗氧量上升，但总效应还是减少心肌耗氧量。

改善心肌缺血区的供血：心率下降，舒张期延长，血液从心外膜流入易缺血的心内膜区。非缺血区与缺血区张力差增加，血液流入缺血区。

抑制脂肪分解酶活性，心肌FFA下降；改善缺血区对葡萄糖的摄取和利用，改善糖代谢，耗氧下降；也能增加组织供氧。
\subsubsection{临床应用}
1、适用于对硝酸酯类不敏感或疗效差的稳定型心绞痛，可减少发作次数，对伴有高血压或心律失常者更为适用；2、缩小心肌梗死范围；3、与硝酸酯类合用互相弥补对抗，耗氧下降

严禁用于冠状动脉痉挛引起的变异性心绞痛，\(\beta\)受体阻断会导致\(\alpha\)受体占据优势，冠脉血管收缩。
\section{钙通道阻滞药}
钙拮抗药——变异型心绞痛
\subsection{硝苯地平、维拉帕米、地尔硫卓}
\subsubsection{作用机制}
1、降低心肌耗氧量：血管平滑肌细胞Ca2+内流减少，心率下降收缩力下降，平滑肌松弛，血压下降。

2、舒张冠状血管：解痉作用，增加灌注，开放侧支循环。

3、保护心肌细胞，防止钙离子超负荷，避免心肌坏死。

4、抑制Ca2+内流，防止血小板聚集
\subsubsection{临床应用}
各种类型心绞痛。

优点：1、适合心肌缺血伴有支气管哮喘者，钙通道阻滞药可以松弛支气管平滑肌；2、强大的扩张冠状动脉作用， 适合变异型心绞痛；3、抑制心肌作用较弱，不易诱发心衰；4、心肌缺血伴外周血管痉挛性疾病患者禁用\(\beta\)受体阻断药，而钙拮抗药因扩张外周血管恰好可用。
\section{其他药物}
卡维地洛：1、阻断\(\beta_{1}\)、\(\beta_{2}\)、\(\alpha\)受体；2、具有一定的抗氧化作用；3、可用于心绞痛、慢性心功能不全和高血压。

尼可地尔：1、激活K+通道，促进外流，细胞膜超极化，抑制Ca2+内流；2、释放NO，增加血管平滑肌细胞内cGMP生成；3、临床变异型心绞痛和慢性稳定型心绞痛，不易产生耐受。
\chapter{作用于血液及造血器官的药物}
[基本要求]

1、了解血液凝固的过程、贫血的分类及形成原因；

2、了解造血细胞生长因子以及血容量扩充药的药理作用和临床应用；

3、掌握抗血小板药、促凝血药、抗贫血药分类、代表药、作用机制和临床应用；

4、掌握肝素和口服抗凝药的作用机理、特点、临床应用与不良反应。
\section{抗凝血药}
凝血有三种通路：

内源性通路：完全靠血浆内的凝血因子逐步使因子X激活从而凝血的通路

外源性通路：被损伤的血管外组织通过释放因子III所引起的凝血通路

共同通路：从受内或外通路激活的因子X开始，到纤维蛋白形成的过程

\subsection{凝血酶间接抑制药}
\subsubsection{肝素}

来源和化学：最初来源于肝，是一种多糖，平均分子量12kDa。

药理作用：抗凝作用，可灭活多种凝血因子。
抗凝作用主要依赖于抗凝血酶III（AT-III），AT-III是凝血酶和因子\( \textup{XII}_{a}\) \(\textup{XI}_{a}\) \(\textup{IX}_{a}\) \(\textup{X}_{a} \)
等含丝氨酸残基蛋白酶的抑制剂。
它与凝血酶通过精氨酸-丝氨酸肽键相结合，形成AT-III-凝血酶复合物而使酶灭活，肝素可加速这一反应达千倍以上。
肝素与AT-III结合可改变其构型，充分暴露活性部位，肝素需要同时与AT-III和因子结合，而低分子量肝素不需要，
复合物形成后肝素则与复合物解离再与其他AT-III结合反复利用。

肝素还具有以下作用：
1、使血管内皮细胞释放脂蛋白酯酶，水解血液中乳糜微粒和VLDL发挥调节血脂作用。
2、抑制炎症介质活性和炎症细胞活动，呈现抗炎作用
3、抑制血管平滑肌细胞增值，抗血管内膜增生
4、抑制血小板聚集

临床应用：
1、血栓栓塞性疾病
2、弥散性血管内凝血
3、防治心肌梗死、脑梗死、心血管手术及外周静脉术后血栓形成
4、体外抗凝

不良反应：
1、出血，主要不良反应，表现为粘膜出血、关节腔积血和伤口出血。可缓慢静脉注射硫酸鱼精蛋白来解救，剂量不可超过50mg。
2、血小板减少症
3、其他，偶有过敏反应。长期使用可致骨质疏松和骨折，孕妇使用可致早产及死胎

禁忌症：对肝素过敏、有出血倾向、血友病、紫癜、严重高血压、细菌性心内膜炎、肝肾功能不全、
溃疡病、颅内出血、活动性肺结核、孕妇、先兆流产、产后、内脏肿瘤、外伤及术后禁用。

相互作用：1、酸性药物不能和碱性药物合用。
2、与阿司匹林等非甾体抗炎药合用可增加出血危险。
3、与糖皮质激素、依他尼酸合用可致胃肠道出血。
4、与胰岛素或磺酰脲类合用可致低血糖。
5、肝素和硝酸甘油可降低肝素活性。
6、与血管紧张素合酶抑制剂合用可引起高血钾。

\subsubsection{低分子量肝素}
分子量低于6.5kDa的肝素，链短，只能与AT-III结合，因此具有选择性抗\( \textup{X}_{a} \)活性。
分子量越小，活性越高。\( t_{\frac{1}{2}} \)较长，每日皮下注射一次即可。

优点：1、抗凝剂量易掌握；2、一般不需要实验室监测抗凝活性；3、毒性小，安全；4、作用时间长；5、可用于门诊患者。

\subsubsection{依诺肝素}
主要用于对抗深静脉血栓，手术后防止静脉血栓，血液透析时防止凝血。
\subsection{凝血酶抑制药}
\subsubsection{阿加曲班}
精氨酸衍生物，与凝血酶催化部位结合，抑制凝血酶的蛋白水解作用，阻碍纤维蛋白原的裂解和纤维蛋白凝块的形成，使某些凝血因子不活化，
抑制凝血酶诱导的血小板聚集及分泌作用，最终抑制纤维蛋白的交联，并促使纤维蛋白溶解。\(t_{\frac{1}{2}}\)短安全范围窄，过量无对抗剂

\subsubsection{水蛭素}
蛋白质，与凝血酶催化部位和阴离子外位点结合抑制凝血酶活性，其余与阿加曲班相同。

\subsubsection{香豆素类}
含有4-羟基香豆素基本结构的物质。常用双香豆素，华法林和醋硝香豆素（新抗凝）。

香豆素类是维生素K拮抗剂，抑制维生素K在肝内由环氧化物向氢醌型转化，阻止维生素K的反复利用。
从而阻止需要维生素K为辅酶的凝血因子XII XI IX X的前体、抗凝血蛋白C和抗凝血蛋白S的\(\gamma\)-羧化作用，使这些因子停留在无活性的前体阶段。
因此香豆素类药物体外起效较慢，体内需要等凝血因子耗尽后才起效，因此起效较慢。

华法林口服吸收快而完全，血浆蛋白结合率高，主要以代谢物排出。
双香豆素口服吸收慢而不规则，与血浆蛋白结合率高药物同服可增加游离药物浓度，增加抗凝效果，经肝药酶羟化失火后由肾排出。
醋硝香豆素主要由原形从肾排出。

主要用于治疗防治血栓栓塞性疾病，心脏瓣膜修复手术后患者需长期服用华法林。

过量易诱发出血，严重可诱发颅内出血，能透过胎盘屏障引起出血，且影响婴儿\(\gamma\)-羧化作用影响胎儿骨骼发育，故孕妇禁用。
应用期间需测定凝血酶原时间，控制在18-24秒为宜。过量用静脉注射大量维生素K或输新鲜血。

不能与阿司匹林、保泰松合用，会使香豆素游离浓度上升。降低维生素K生物利用度的药物或广谱抗生素均能增强香豆素类的作用，肝病也可增强作用。
肝药酶诱导剂苯巴比妥、苯妥英钠、利福平可加速香豆素类代谢，降低抗凝作用。

\section{纤维蛋白溶解药和纤维蛋白溶解抑制药}
\subsection{纤维蛋白溶解药}
能够使纤维蛋白溶解酶原转变为纤维蛋白溶解酶，纤溶酶降解纤维蛋白和纤维蛋白原，限制血栓增大和溶解血栓。又称血栓溶解药。
\subsubsection{链激酶}
第1代，易全身纤溶，出血，分子量47kDa，\( t_{\frac{1}{2}} \)=23min，静脉给药。
与内源性纤溶酶原结合成复合物，促使纤溶酶原转变为纤溶酶，水解血栓中的纤维蛋白，血栓溶解。

主要用于治疗血栓栓塞性疾病。

不良反应:引起出血，注射局部可出现血肿。
\subsubsection{尿激酶}
可被纤溶蛋白酶激活剂的抑制物（PAI）中和，产生的纤溶酶可被\(\alpha\)-抗纤溶酶（\(\alpha\)-AP）灭活，故需要耗竭PAI和\(\alpha\)-AP后才能生效。

尿激酶无抗原性，可用于链激酶过敏者。
\subsubsection{阿尼普酶}
又称茴香酰化纤溶酶原-链激酶激活剂的复合物。活性中心与一个酰基结合而被封闭
进入血液后弥散到血栓含纤维蛋白表面，通过复合物的赖氨酸纤溶酶原活性中心与纤维蛋白结合，缓慢脱掉乙酰基后，变为纤溶酶而发挥溶解血栓作用。
有一定潜伏期，但不影响与纤维蛋白的结合力，因此具有溶栓选择性。

优点有：1、在体内缓慢活化，可静脉注射提高结合量，在血中不受\( \alpha _{2}\)-抗纤溶酶的抑制；
2、本品与赖-纤溶酶原的复合物较易进入血凝块与纤维蛋白结合，而谷-纤溶蛋白要降解为与赖-纤溶酶原才能结合，因此很少引起全身性反应，出血少。

可用于急性心肌梗死和其他血栓疾病

可导致长时间血液低凝状态，出血常发生在注射部位和胃肠道，类似链激酶有抗原性
\subsubsection{阿替普酶}
t-PA，t-PA在靠近纤维蛋白-纤维蛋白溶酶原相结合的部位，通过赖氨酸残基与纤维蛋白结合，激活与纤维蛋白结合内源性纤溶酶原转变为纤溶酶。主要用于心梗患者

优点：1、溶栓效率高，生效快，耐受性好；2、生产成本低，给药方法简便，不需要按体重调整剂量。

不良反应：出血、血小板减少、出血倾向慎用。
\subsection{纤维蛋白溶解抑制药}
\subsubsection{氨甲苯酸}
竞争性抑制纤溶酶原激活因子，使其不能转变为纤溶酶。主要用于纤维蛋白溶解症导致的出血，对癌症出血、创伤出血等非纤维蛋白溶解引起的出血无效果。
\section{抗血小板药}
\subsection{抑制血小板花生四烯酸代谢}
\subsubsection{阿司匹林}
环氧酶抑制药抑制COX-1，抑制花生四烯酸转变为\(\textup{PGG}_{2}\)和\(\textup{PGH}_{2}\)，从而使血小板\(\textup{TXA}_{2}\)合成减少。

1、预防：小剂量用于预防慢性稳定性心绞痛，心肌梗死的一级和二级预防，脑梗死、脑卒中或短暂性缺血发作后脑梗死的二级预防，预防瓣膜修复术后或冠脉搭桥术后的血栓；
2、治疗：冠状动脉硬化性心脏病、心肌梗死和缺血性脑血管病。能减少缺血性心脏病发作和复发的风险，也可以使一过性脑缺血发作患者的卒中发生率死亡率降低。
\subsubsection{利多格雷}
\(\textup{TXA}_{2}\)合酶抑制药，并且能够中等程度的拮抗\(\textup{TXA}_{2}\)受体。对血小板血栓和冠状动脉血栓的作用比阿司匹林有效。
\subsection{增加血小板内c-AMP}
\subsubsection{双嘧达莫}
抑制血小板聚集，体内外均抗血栓：
1、抑制磷酸二酯酶活性，细胞内cAMP上升，抑制血小板聚集；2、增强血管内皮细胞PGI2生成和活性；3、激活腺苷酸环化酶， cAMP上升；4、轻度抑制血小板的环氧酶 TXA2合成

主要用于防治血栓栓塞性疾病，防止血小板血栓形成；抑制动脉粥样硬化早期病变。
\subsection{抑制ADP活化血小板}
\subsubsection{噻氯匹定}
ticlopidine ----抑制血小板活化和粘附：

抑制ADP诱导的\(\alpha\)颗粒分泌；
抑制ADP诱导的血小板膜糖蛋白受体（\(\textup{GPII}_{b}\)/\(\textup{III}_{a}\)）与纤维蛋白原结合位点的暴露；
拮抗ATP对腺苷酸环化酶的抑制。

预防脑中风，心肌梗死及外周动脉血栓性疾病的复发，疗效优于阿司匹林

不良反应：中性粒细胞减少（2.4％）（用药3个月内定期检查血象），恶心、腹泻
\subsection{\(GPII_{b}/III_{a}\)受体阻断药}
血小板膜糖蛋白Ⅱb/Ⅲa受体（GP Ⅱb/Ⅲa R）是ADP、凝血酶、TXA2等致血小板聚集的最终共同通路。

\(\textup{GPII}_{b}\)/\(\textup{III}_{a}\)R配体：纤维蛋白原、血管性血友病因子、内皮诱导因子等；
\(\textup{GPII}_{b}\)/\(\textup{III}_{a}\)受体阻断药阻碍血小板同上述配体联结，从而抑制血小板聚集；
\(\textup{GPII}_{b}\)/\(\textup{III}_{a}\)受体单克隆抗体：Abciximab（阿昔单抗），特异性优于aspirin；

用于急性心梗、溶栓、不稳定型心绞痛和血管成形术后再梗死的治疗，效果较好。抑制血小板聚集作用强，不良反应少。

\section{促凝血药}
\subsubsection{维生素K}
\(\gamma\)-羧化酶的辅酶，参与肝合成凝血因子XII XI IX X、抗凝血蛋白C、抗凝血蛋白S等的活化过程。缺乏维生素K肝脏仅能合成无凝血活性的。

主要用于梗阻性黄疸、胆瘘、慢性腹泻、早产儿、新生儿出血等患者和香豆素、水杨酸导致凝血酶原过低的出血者，也可用于广谱抗生素导致的维生素K缺乏患者
\subsubsection{凝血酶}
从猪、牛血液中提取的无菌制剂，直接作用于纤维蛋白原。通常用于止血困难的小血管、毛细血管和实质性脏器出血。
\section{抗贫血药及造血细胞生长因子}
贫血：循环血液中红细胞数或血红蛋白长期低于正常值的一种病理现象。

缺铁性贫血（iron deficiency anemia）：血液损失过多或 铁盐吸收不足所致红细胞呈小细胞，低色素性 治疗：补充铁剂

巨幼红细胞性贫血（megaloblastic anemia）：叶酸或VitB12缺乏，红细胞呈大细胞，高色素性，白细胞及血小板降低 治疗：补充叶酸和维生素VB12。

再生障碍性贫血（aplastic anemia）：骨髓造血功能障碍，红细胞、粒细胞、血小板下降，没得治。
\subsection{抗贫血药}
\subsubsection{铁剂}
主要在十二指肠及空肠上端吸收，无机铁以\(\textup{Fe}^{2+}\)的形式吸收，\(\textup{Fe}^{3+}\)难以吸收，络合铁更易吸收。
吸收到骨髓的铁，吸附在有核红细胞膜上并进入细胞内的线粒体，与原卟啉结合形成血红素，再与珠蛋白结合形成血红蛋白。
促进吸收：胃酸、VitC、还原性物质；妨碍吸收：胃酸减少、高磷、高钙、鞣酸、 四环素。
可氧化成Fe3+与去铁蛋白结合成铁蛋白储存或由转铁蛋白转运至肝、脾、骨供造血或储存。由肠粘膜上皮细胞、胆汁、尿、汗排泄。

铁剂会刺激胃肠道，有时会引起便秘。小儿误服1g以上铁剂会急性中毒，表现为坏死性肠胃炎，以磷酸盐或碳酸盐洗胃并用去铁胺注入胃中结合残存的铁。
\subsubsection{叶酸}
叶酸在体内被还原和甲基化为(5-CH3-FH4), (5-CH3-FH4)作为甲基供体使VitB12转成甲基B12，自身变为四氢叶酸（FH4）

FH4作为一碳基团转移酶的辅酶，1、传递一碳基团：参与嘌呤核苷酸合成；2、dUMP合成胸腺嘧啶脱氧核苷酸（dTMP）；3、氨基酸互变：丝氨酸\(\rightarrow\)甘氨酸。

叶酸缺乏会导致核苷酸特别dTMP合成受阻\(\rightarrow\)有丝分裂障碍\(\rightarrow\)巨幼红细胞性贫血

用于治疗各种巨幼红细胞性贫血，但甲氨蝶呤等抗叶酸药引起的巨幼红细胞性贫血，给叶酸无效因为二氢叶酸还原酶被抑制，需要使用甲酰四氢叶酸钙。
叶酸与VB12一起使用效果更好
\subsubsection{VB12}存在形式：氰钴胺，羟钴胺，甲钴胺，5'-脱氧腺苷钴胺。
必须与“内因子”结合\(\rightarrow\)回肠特异转运系统\(\rightarrow\)吸收\(\rightarrow\)储存肝脏和参与生化反应

药理作用：1、促进四氢叶酸循环利用，VitB12起传递甲基；2、脱氧腺苷B12甲基丙二酰CoA转变为琥珀酰CoA的辅助因子，缺乏会导致神经髓鞘脂质合成障碍从而导致神经炎

临床：恶性贫血、巨幼红细胞性贫血和神经系统疾病辅助治疗。可致过敏反应，不能滥用，不可静脉给药

\subsection{造血细胞生长因子}
\subsubsection{红细胞生成素}
由肾皮质近曲小管管周细胞分泌的一种糖蛋白；临床应用的是重组人红细胞生成素促红素；

药理作用：与红系干细胞表面EPO受体结合，导致细胞内磷酸化及钙离子上升，导致红系干细胞增生和成熟，并促使网织红细胞从骨髓中释放入血。

用法：静注或皮下注射，治疗慢性肾性贫血和化疗、艾滋病药物引起的贫血。

不良反应：血压升高；血凝增强 ；局部皮肤反应及关节痛。
\subsubsection{粒细胞集落刺激因子}
重组人G-CSF非格司亭，由血管内皮细胞、单核细胞和成纤维细胞合成的糖蛋白，临床使用非格司亭；

药理作用：刺激粒细胞集落形成单位；促使中性粒细胞成熟；刺激粒细胞从骨髓释出，增强趋化和吞噬功能；对巨噬、巨核细胞影响小；

临床应用：用于骨髓移植、化疗等引起的粒细胞缺乏症

不良反应：可产生轻、中度骨痛，及局部反应。
\subsubsection{粒细胞/巨噬细胞集落刺激因子}
沙格司亭—从酵母菌产生的重组人GM-CSF；

药理作用：刺激造血前体细胞增殖、分化；刺激粒细胞、单核细胞、T细胞生长，诱导生成粒细胞、巨噬细胞集落形成单位及粒细胞/巨噬细胞集落形成单位；
对红细胞增生也有影响；促进巨噬细胞和单核细胞对肿瘤细胞的裂解。

临床应用：用于骨髓移植、化疗、再障或艾滋病有关的白细胞或粒细胞缺乏症；

不良反应：引起骨痛、流感样症状、发热、腹泻及呼吸困难等不良反应。

\section{血容量扩充药}
\subsubsection{右旋糖酐}
葡萄糖聚合物，分为不同分子量。主要用中、低、小分子量。主要用于提高血浆胶体渗透压，扩充血容量，维持血压。作用强度与维持时间随分子量降低而降低。
低、小分子量右旋糖酐有阻止红细胞血小板及纤维蛋白聚合，降低血液黏性，并且能抑制凝血因子II，改善微循环。

用于低血容量休克，低小分子量可用于中毒性外伤性和失血性休克，可防止休克后期DIC。

偶见过敏反应，少量大分子右旋糖酐蓄积可导致凝血障碍和出血。

\chapter{影响自体活性物质的药物}
掌握抗组胺药的药理作用、临床应用。

自体活性物质（autacoids)：以旁分泌的形式达到邻近部位，而不进入血液循环，又称局部激素。

如：PG、5-HT（5-羟色胺）、组胺、白三烯、血管活性肽类（P物质，激肽类，血管紧张素，利尿钠肽等）、NO、腺苷等。

特点：是由组织而非特定内分泌腺、产生，不需由血循环运送的靶器官发挥作用。
\section{组胺}
临床用判断有无真性胃酸缺乏症，现在多用五肽促胃酸激素代替。

培他司汀 betahistine，抗眩啶， H1受体激动剂，促进脑干和迷路血循环，纠正内耳血管痉挛，减轻膜迷路水肿，抗血小板聚集和抗血栓形成。用于内耳旋晕病，头痛，慢性缺血性脑血管病。

英普咪定 impromidine，甲双咪胍 选择性H2受体激动剂，胃功能检查。

受体有4种亚型：\(\textup{H}_{1}\)、\(\textup{H}_{2}\)、\(\textup{H}_{3}\)和\(\textup{H}_{4}\)
\subsubsection{\(\textup{H}_{1}\)R效应}
兴奋支气管和胃肠道平滑肌；部分扩张血管，增加毛细血管通透性等。H1R激动剂：2-甲基组胺
\subsubsection{\(\textup{H}_{2}\)R效应}
H2R效应：刺激胃液分泌，扩张血管，兴奋人心脏（正性肌力）。H2R激动剂：英普咪定
\subsubsection{\(\textup{H}_{3}\)R效应}
存在于中枢和外周神经末稍，抑制谷氨酸释放，使自身释放减少。H3R激动剂：（R）\(\alpha\)-甲基组胺
效应：

\subsubsection{心脏}
心肌收缩力上升，豚鼠则下降
\subsubsection{血管}
H1R、H2R激动\(\rightarrow\)小A、V扩张\(\rightarrow\)回心血量\(\downarrow\)；
激动 H1R\(\rightarrow\)毛细血管扩张\(\rightarrow\)通透性\(\uparrow\)\(\rightarrow\)局部水肿和全身血液浓缩；
人的冠脉有H1、H2受体，二者失衡可致冠脉痉挛。
皮内注射大剂量组胺：强而持久降压，甚至休克；

注射小剂量组胺“三重反应”：1、毛细血管扩张红斑；2、血管通透性增加，红斑上形成丘疹；3、轴索反射致小动脉扩张， 丘疹周围形成红晕；

麻风患者“三重反应”常不完全：因为皮肤神经受损
\subsubsection{血小板}
存在H1、H2R、H1R激动，激活PLA2，介导AA释放，调节细胞内Ca2+水平，促进血小板聚集。另一方面激动H2R可增加血小板中的cAMP，对抗血小板聚集，最终结果视二者的功能平衡变化而定。
\subsubsection{腺体}
激动胃壁细胞H2R，胃酸分泌\(\uparrow\)，较弱的唾液肠液和支气管腺体分泌\(\uparrow\)
\subsubsection{平滑肌}
支气管：健康人常亮无影响；哮喘患者敏感，少量即可引起呼吸困难；豚鼠最敏感。

胃肠道平滑肌：多种动物收缩，豚鼠回肠最敏感，人大量可引起收缩。

子宫：大鼠松弛；豚鼠收缩；人无影响。

\section{抗组胺药}
\subsection{\(\textup{H}_{1}\)受体阻断剂}
\subsubsection{第一代}
苯海拉明diphenhydramine，苯那君

异丙嗪 promethazine，非那根

曲吡那敏 pyribenzamine，扑敏宁

氯苯那敏 chlorpheniramine，扑尔敏
\subsubsection{第二代}
西替利嗪（记住） cetirizine，仙特敏

美喹他嗪 mequitazine，甲喹酚嗪

阿司咪唑 astemizole，息斯敏； 阿伐斯汀

acrivastine，新敏乐； 氯雷他定（记住）
\subsubsection{药理作用}
口服注射均易吸收，可完全对抗组胺收缩胃肠和支气管平滑肌的作用；大部分对抗其扩血管及增加毛细血管通透性的作用；仅部分对抗组胺所致的降压和心脏作用。

第一代有中枢抑制作用因此会导致嗜睡，第二代无中枢作用

苯海拉明等具有抗胆碱作用，止吐、防晕动作用较强，咪唑斯汀治疗鼻塞效果明显。

\subsubsection{临床应用}
1、皮肤粘膜变态反应性疾病：荨麻疹、过敏性鼻炎的首选药（重点），昆虫咬伤所致瘙痒和水肿较好，接触性皮炎、血清病有效，支气管哮喘效差，过敏性休克无效（重点）；
2、晕动病、放射病所致呕吐；
3、与其它药物配伍应用。
\subsubsection{不良反应}
1、CNS反应：第一代多见镇静、嗜睡、乏力。驾驶员和高空作业者工作期间不宜使用；
2、消化道反应：口干、厌食、腹泻、呕吐、便秘等消化道症状；
3、其它反应：偶见粒细胞减少和溶血性贫血。

\subsection{\(\textup{H}_{2}\)受体阻断剂}
西咪替丁cimetidine，甲氰咪胍

雷尼替丁 ranitidine，甲硝呋胍

法莫替丁 famotidine，

尼扎替丁 nizatidine，

罗沙替丁 roxatidine

〔体内过程〕
口服吸收迅速，西咪替丁能通过血脑屏障和胎盘屏障。

\subsubsection{药理作用}
1、抑制胃酸分泌：与组胺可逆性竞争H2受体。减少胃酸分泌，抑制基础胃酸和食物、组胺、五肽胃泌素、Ach所致的胃酸分泌，作用强，胃酸和胃蛋白酶均下降，促进溃疡愈合。法莫替丁作用最强>雷尼替丁>西咪替丁

2、心血管系统：可拮抗正肌、正频。部分拮抗扩血管、降压作用。在抑制胃酸分泌的剂量对心血管影响小。

3、免疫功能调节：逆转组胺对免疫系统有抑制作用，细胞免疫和体液免疫功能均有所降低。
\subsubsection{临床应用}
1、胃和十二指肠溃疡，减轻疼痛，促进愈合。

2、各种原因引起的免疫功能低下和肿瘤的辅助治疗。

易于耐受，不良反应少。
\subsection{\(\textup{H}_{3}\) \(\textup{H}_{4}\)受体阻断剂}
\subsubsection{H3受体}
阿尔茨海默病、注意力缺陷多动症、帕金森病等神经行为失调有关， H3受体拮抗剂能改善大鼠的学习与记忆能力，其拮抗剂如GT2277等的应用正在临床试验中。
\subsubsection{H4受体}
主要在与炎症反应有关的组织和造血细胞中表达。可能是一种重要的炎症性受体。
\chapter{作用于呼吸系统的药物}
了解祛痰药的分类；掌握镇咳药的分类及应用；重点掌握平喘药的分类、作用与用途
\section{平喘药}

\subsection{抗炎平喘药}
\subsubsection{糖皮质激素}
\subsubsection{药理作用}
GCs进入靶细胞内与糖皮质激素受体结合成复合物，再进入细胞核内调控炎症相关靶基因的转录，通过抑制哮喘时炎症反应的多个阶段发挥平喘作用。

1、抑制多种参与哮喘发病的炎性细胞和免疫细胞功能；2、抑制细胞因子和炎症介质的产生；3、抑制气道高反应性；4、增强支气管和血管平滑肌对儿茶酚胺的敏感性。
\subsubsection{临床应用}主要以气雾吸入方式在呼吸道局部给药。

\subsubsection{磷酸二酯酶-4抑制剂：罗氟司特}
\subsubsection{药理作用}PDE-4主要分布于炎症细胞、气泡上皮细胞和平滑肌细胞内，PDE-4是细胞内特异性的cAMP水解酶，PDE-4抑制剂抑制PDE-4活性，增高细胞内cAMP而起效

1、抑制炎症细胞聚集和活化；2、扩张气道平滑肌，轻度，缓解气道高反应性；3、缓解气道重塑，减少上皮细胞基底的胶原蛋白沉着、气道平滑肌增厚、杯状细胞增生和黏蛋白的分泌；促进气道上皮纤毛运动而促进排痰。
\subsubsection{临床应用}因糖皮质激素治疗COPD不能明显改善肺功能，也无法降低死亡率，所以罗氟司特被用于治疗反复发作并加重的成人重症COPD，常与长效支气管扩张药联合应用。
\subsection{支气管扩张药}
\subsubsection{肾上腺素受体激动药\(\beta_{2}\)R为主}
\subsubsection{药理作用}
气道平滑肌松弛，抑制肥大细胞与中性粒细胞释放炎症介质，增强气道纤毛运动，促进黏液分解，降低血管通透性，减轻气道黏膜下水肿

非选择性肾上腺素受体激动药：肾上腺素，用于哮喘急性发作；异丙肾，可产生严重不良反应；麻黄碱，缓慢、温和、持久，可口服，用于轻症哮喘或预防哮喘发作

选择性（\(\beta_{2}\)受体激动药：沙丁胺醇：扩张支气管作用与异丙肾上腺素相近；克仑特罗：中效\(\beta_{2}\)R激动药 瘦肉精

\subsubsection{茶碱类}
\subsubsection{药理作用}
1、抑制磷酸二酯酶，非选择性PDE抑制剂，能降低细胞内cAMP水平，舒张支气管平滑肌

2、阻断腺苷受体，减轻内源性腺苷受体引起的气道收缩效应

3、增加内源性儿茶酚胺的释放

4、免疫调节与抗炎作用，抑制肥大细胞、嗜酸性粒细胞、巨噬细胞、T淋巴细胞的功能，减少炎症介质释放，降低微血管通透性，降低气道炎症反应。

5、增加膈肌收缩力并促使支气管纤毛运动。
\subsubsection{临床应用}
1、支气管哮喘。对于\(\beta\)受体激动剂不能控制的急性哮喘可以静脉注射茶碱。慢性哮喘患者可以口服氨茶碱。

2、慢性阻塞性肺病。对COPD伴有喘息、COPD伴有新功能不全的心源性哮喘患者有明显的疗效。因为茶碱还有扩张肺动脉及降低肺动脉压、强心、利尿作用。

3、中枢型睡眠呼吸暂停综合征。中枢兴奋作用，对脑部疾病或原发性呼吸中枢病变导致通气不足患者有效。
\subsubsection{不良反应}
治疗窗口较窄，主要不良反应为胃肠道不良反应、中枢兴奋、急性中毒：心动过速、心律失常等。

\subsubsection{抗胆碱药}
\subsubsection{异丙托溴铵}
对\(\beta_{2}\)R激动药耐受的患者有效。
\subsection{抗过敏平喘药}
\subsubsection{炎症细胞膜稳定剂}
\subsubsection{色甘氨酸钠}
1、稳定肥大细胞膜；2、抑制气道感觉神经末梢功能与气道神经源性炎症；3、阻断炎症细胞介导的反应。

预防用药物。
\subsubsection{H1受体阻断剂}
\subsubsection{酮替芬}
除类似色甘氨酸钠以外还有有H1受体阻断作用。能预防和逆转\(\beta_{2}\)R的“向下调节”，加强其平喘作用。
\subsubsection{白三烯阻断药}白三烯是炎症介质。
\section{镇咳药}
中枢型：延髓咳嗽中枢的选择性抑制作用。有磷酸可待因、氢溴酸右美沙芬。可待因有成瘾性。

外周性：盐酸那可汀。可抑制肺牵张反射引起的咳嗽，兼具兴奋呼吸中枢。
\section{祛痰药}
痰液稀释剂：氯化铵、愈创木酚甘油醚。

黏痰溶解药：乙酰半胱氨酸、脱氧核糖核酸酶。

黏痰调节药：溴己新
\chapter{作用于消化系统的药物}
了解消化系统药物的分类

了解泻药、止泻药和利胆药

掌握止吐药、增强胃肠动力药的分类及应用

重点掌握抗溃疡病药的分类及其作用机理和主要代表药物
\section{治疗消化性溃疡药}
消化性溃疡，常见病。由于胃及十二指肠多发，又称胃及十二指肠溃疡。

溃疡发生的因素：遗传；地理位置；饮食；情绪（长期精神过度紧张，焦虑不安）；烟酒；药物。

直接原因：胃酸、胃蛋白酶、幽门螺杆菌等攻击因子超过胃黏液、HCO分泌和胃粘膜等防御因子的承受能力

主要作用：胃酸浓度下降、胃酶活性下降、攻击因子作用减弱。胃黏液、\(\textup{HCO}_{3}^{-}\)分泌上升，胃肠粘膜保护作用增强，防御因子作用增强。
\subsection{抗酸药}
能中和胃酸的弱碱性无机药物。作用持久，不吸收，不产气，不致腹泻和便秘，对胃粘膜及溃疡面有保护和收敛作用。

氢氧化铝、氢氧化镁、三硅酸镁、碳酸钙、碳酸氢钠。

仅能中和已经分泌的胃酸，不能调节分泌，分首选药。
\subsection{抑制胃酸分泌药}
\subsubsection{H2受体阻断药}
H2R兴奋，细胞内cAMP上升，质子泵兴奋，促进胃酸分泌。

阻断H2R可：1、基础胃酸，夜间胃酸分泌均下降（重点）；2、抑制胃泌素、进食、迷走神经兴奋引起的胃酸分泌；3、胃蛋白酶分泌也下降。入睡前服用是治疗十二指肠溃疡的首选。

西咪替丁：消化性溃疡，上消化道出血，促胃液素分泌瘤。不良反应主要为：CNS：焦虑、定向障碍、幻觉；内分泌：抗雄激素、促催乳素分泌；其他：心动过缓、肝肾损害；肝药酶抑制作用。
\subsubsection{H+-K+-ATP酶（质子泵）抑制药}
多种胃酸分泌刺激的效应环节，影响该酶的药物作用强大。（重点）

奥美拉唑：1、能选择性的浓集于parietal Cell膜分泌小管系统与壁细胞H+-K+-ATP酶不可逆结合并灭活该酶，抑制胃酸胃蛋白酶分泌。
2、对胃泌素、组胺、Ach、食物等因素诱发的均可抑制。24h抑制率达95；
3、清晨口服20 mg可控制胃内pH在3以上达16-18 h；
4、大剂量可达无酸状态；
5、有抗幽门螺旋杆菌作用。

临床应用：良性胃及十二指肠溃疡；术后溃疡、应激性溃疡；急性胃粘膜出血、促胃液素瘤；反流性食道炎。

不良反应：皮疹、外周神经炎、性激素紊乱；口干、恶心、呕吐、腹胀。

奥美拉唑：一代；兰索拉唑：二代；半妥拉唑、雷贝拉唑：三代。泰妥拉唑：新型。
\subsubsection{M受体阻断药}
哌仑西平：药理作用：1、抑制\(\textup{M}_{3}\)受体，对\(\textup{M}_{3}\)受体选择性低；2、抑制胃酸分泌；3、阻断Ach对胃粘膜中ECL细胞的激动作用，减少His释放；4、抑制G细胞MR激动作用，减少胃泌素释放；5、解痉作用。

目前已少用于溃疡治疗
\subsubsection{胃泌素受体阻断药}
丙谷胺：1、抑制胃酸分泌；2、促进胃黏膜粘液合成；3、增强粘液-碳酸氢盐屏障。

用于胃、十二指肠溃疡和胃炎。疗效差，很少单独使用。偶尔有口干、失眠等症状。
\subsection{增强胃粘膜屏障功能的药物}
胃粘膜屏障:细胞屏障和其表面粘液-碳酸氢盐屏障，能防止有害因子损伤胃粘膜，粘膜上  皮能迅速重建和再生，有修复作用。屏障功能受损可导致溃疡。

米索前列醇：与胃壁细胞和胃粘膜的前列腺素受体结合，可刺激胃粘液和HCO3-分泌，维持粘膜细胞完整性。

①抑制基础胃酸分泌；②抑制组胺、胃泌素，食物刺激所致胃酸分泌；③胃蛋白酶分泌减少。

临床应用：胃、十二指肠溃疡、胃炎所致出血；非甾体抗炎药所致胃粘膜损伤、溃疡，防止复发。

不良反应：轻；兴奋子宫，致流产，孕妇禁用。 

硫糖铝：1、与溃疡面的亲和力是正常粘膜的6倍，形成保护屏障；2、促进胃和十二指肠粘膜合成PGE2；3、增强表皮生长因子、碱性成纤维细胞生长因子作用，使之聚集于溃疡区，促进溃疡愈合；4、抑制幽门螺杆菌繁殖，阻止该菌的蛋白酶和脂酶对粘膜的破坏。

注意：1、在酸性环境中发挥作用，不宜与碱性药合用；与布洛芬、吲哚美锌、氨茶碱、四环素、地高辛合用，能降低上述药物的生物利用度；3、减少甲状腺素的吸收。

麦滋林：由99的谷氨酰胺和0.3的水溶性薁组成，1、前者增加PGE2合成，粘液合成，促进粘膜细胞增殖；2、后者抗炎，抑制胃蛋白酶活性。

\subsection{抗幽门螺旋杆菌药}
G-厌氧菌   慢性胃窦炎
清除幽门螺杆菌可明显降低十二指肠溃疡复发率。

甲硝唑、四环素、阿莫西林、罗红霉素为常用药。

单一应用易产生耐药性，2-3种联合用药。

质子泵抑制剂或铋制剂:：美拉唑 40mg/d、兰索拉唑 60mg/d、枸橼酸铋钾480 mg/d。
抗菌药物：克拉霉素500-1000mg/d、阿莫西林1500-2000mg/d、甲硝唑 800mg/d。

\section{消化功能调节药}
\subsection{助消化药}
用于消化不良。为消化液成分或为促进消化液分泌药：
稀盐酸；胃蛋白酶；胰酶；

干酵母：富含多种B族维生素

乳酶生：分解糖类，产生乳酸
\subsection{止吐药}
H1 受体阻断药：如苯海拉明，中枢镇静和止吐，防治晕动病。

M 受体阻断药：东莨菪碱、阿托品。阻断呕吐中枢和外周反射途径中MR，降低迷路感受器的敏感性和抑制小脑前庭通路的传导，抗晕动病。

多巴胺（D2）受体阻断药：

甲氧氯普胺metoclopramide（灭吐灵）：

CNS：作用于中枢化学感受区（CTZ），阻断D2 受体，高剂量阻断5-HT3 受体；外周：胃肠道，促进运动，增加贲门括约肌张力，松弛幽门，加速胃排空和肠内容物运动。

应用：慢性功能性消化不良引起的胃肠运动障碍。

不良反应：锥体外系症状。

多潘立酮（吗丁啉）：1、阻断CTZ和上消化道的D2受体；2、具有胃肠推动和抗呕吐作用；3、不易通过血脑屏障，故锥体外系反应罕见（重点）。

用于治疗慢性食后消化不良、恶心、呕吐和胃潴留，对中度化疗药引起的呕吐有效。

5-HT3 受体拮抗药：昂丹司琼：

竞争性拮抗CNS和胃肠道5-HT3受体；明显抑制抗癌药顺铂所致的呕吐；疗效优于甲氧氯普胺。
\subsection{增强胃动力药}
MR激动药和胆碱酯酶抑制剂；D2R阻断药；5-HT4R激动药：西沙比利
\subsection{止泻药与吸附药}
阿片制剂：非细菌感染性腹泻；

地芬诺酯：合成哌替啶衍生物，主作用于胃肠μ-R，CNS副作用小。用于功能性腹泻；

鞣酸蛋白：凝固、沉淀肠粘膜表面蛋白质，收敛作用；

药用炭：吸附剂，起止泻和阻止毒物吸收作用
\subsection{泻药}
1、刺激性：酚肽，蒽醌类；用于急、慢性便秘，腹部X线、内窥镜检查及术前排空 肠内容物等。

2、渗透性：硫酸镁，乳果糖；主要用于排除肠内毒物及某些驱肠虫药服后连虫带药一起排出

3、润滑性：液体石腊。适用于老人、痔疮及肛门手术等患者。
\subsection{利胆药}促进胆汁分泌或胆囊排空药物；常涉及胆汁酸。一般用于胆石症、胆囊炎的治疗。
\section{思考题}
抗消化性溃疡药物的分类及代表药；

抑制胃酸分泌的药物分类，作用机制，代表药物；

止吐药的分类、机制和代表药物。
\chapter{肾上腺皮质激素类药物}
了解肾上腺皮质激素的分类、构效关系，掌握肾上腺皮质激素的作用、用途与不良反应。
主要有糖皮质激素、盐皮质激素
\subsection{糖皮质激素}
主要调剂糖、蛋白质、脂肪代谢。对水盐代谢影响小。由胆固醇衍化而来，又称类固醇激素。基本结构是甾核，又称甾体激素。

受刺激--》进入应激性“紧急状态”--》刺激大脑皮层和肾上腺髓质--》肾上腺素增多--》刺激垂体前叶分泌释放促皮质激素--》促使肾上腺皮质分泌糖皮质激素。

主要在肝脏代谢，由尿排出。可的松、泼尼松在肝脏内分别转化为氢化可的松和泼尼松龙发挥作用---严重肝功能不良病人
\subsubsection{对代谢的影响}
1、促进糖原异生、抑制葡萄糖分解、减少机体对葡萄糖的利用。长期大量应用可导致血糖升高甚至发生继发性糖尿病。糖尿病人不用，低糖饮食。

2、促进蛋白质分解代谢、抑制合成。长期大量应用可导致淋巴组织和肌肉组织萎缩、皮肤变薄。骨的蛋白质分解代谢增强导致骨质疏松。高蛋白饮食、必要时合用蛋白质同化激素。

3、促进脂肪分解和脂肪重新分布，可导致向心性肥胖。

4、较弱的盐皮质激素作用，保钠排钾；长期用药骨质脱钙、引起骨质疏松。
\subsubsection{抗炎机理}
炎症早期：提高血管紧张性，使通透性降低；抑制白细胞浸润；减少各种炎症因子释放。缓解炎症引起的红、肿、热、痛等症状。

炎症晚期：抑制毛细血管和纤维母细胞增生，缓解肉芽组织形成，坚强炎症引起的粘连和瘢痕形成，减轻后遗症。
降低机体防御功能: 可致感染扩散，阻碍创口愈合。 但抗炎不抗菌
\subsubsection{特点及利与弊}
\subsubsection{抗内毒素、抗休克机理}
\subsubsection{不良反应与禁忌症}

\end{document}